% ============================================================
% Chapter 1: Introduction
%
% PURPOSE:
%   Establish your problem statement and research questions.
%   Convince the reader that your topic is relevant, important,
%   and worth studying. Do not go into deep technical detail
%   here — save that for later chapters.
%
% SUGGESTED SECTIONS (adapt to your needs):
%   1.1 Background and Context
%   1.2 Problem Statement
%   1.3 Research Questions or Objectives
%   1.4 Significance
%   1.5 Thesis Organization
%
% TIPS:
%   - Start broad, then narrow toward your specific problem.
%   - State your research questions explicitly.
%   - Cite a handful of sources to ground your topic in the
%     literature. Use \cite{KEY} where KEY matches an entry
%     in References.bib. Example: \cite{smith2024}.
%   - Every \cite{} call produces a number like [1] in the text.
%     Multiple citations: \cite{smith2024, jones2023}.
%
% TARGET LENGTH: 4--8 pages.
% ============================================================

\chapter{Introduction}
\label{chapter1}


\section{Background and Context}

\noindent
[Introduce your research area. What field are you working in?
What is the current state of the field, and why does it matter?
Cite a few foundational sources to establish that this is a
recognized area of interest \cite{example_journal}.]


\section{Problem Statement}

\noindent
[State the specific problem your thesis addresses. What gap exists
in current knowledge, practice, or technology? Why has it not been
solved? Be as precise as possible about what you are doing and for
whom \cite{example_conference}.]


\section{Research Questions}
% If your thesis is objective-driven rather than question-driven,
% rename this section "Objectives" and use a numbered list instead.

\noindent
This thesis addresses the following research questions:

\begin{enumerate}
    \item[\textbf{RQ1}] [State your first research question here.]
    \item[\textbf{RQ2}] [State your second research question here.]
    % Add or remove questions as needed.
\end{enumerate}


\section{Significance}

\noindent
[Explain why solving this problem matters. Who benefits? How does
your work build on or differ from what already exists?]


\section{Thesis Organization}

\noindent
[Briefly describe what each chapter covers. One or two sentences
per chapter is sufficient. Example below — update to match your
actual structure.]
\newline

\noindent
The remainder of this thesis is organized as follows.
Chapter~\ref{chapter2} reviews the relevant background and
related work. Chapter~\ref{chapter3} describes the methodology.
The thesis concludes with a summary of contributions and directions
for future work.
\newline
